\documentclass[conference]{IEEEtran}
\IEEEoverridecommandlockouts

\usepackage[spanish,es-nodecimaldot,es-noshorthands]{babel}
\usepackage[utf8]{inputenc}
\usepackage[T1]{fontenc}
\usepackage{hyperref}

\begin{document}

\title{Conclusión Individual: Proyecto Final \\
\large Terminal de Juegos Retro Distribuida con Mando Embebido}

\author{
\IEEEauthorblockN{Dorantes Mejía Emmanuel}
\IEEEauthorblockA{\textit{Ingeniería en Computación — Facultad de Ingeniería, UNAM} \\
No. Cuenta: 313004766 \quad | \quad Grupo 6, Equipo 2}
}

\maketitle

\section*{Aportación Técnica al Sistema}

Durante la realización del proyecto final, mi aportación principal consistió en el diseño e implementación del subsistema de adquisición de señales analógicas en la Raspberry Pi Pico, así como en la definición de las estrategias de sincronización y gestión de memoria estática en el firmware \textit{bare-metal}.

Desarrollé los módulos \texttt{read\_adc.c} y \texttt{read\_adc.h} utilizando las bibliotecas de bajo nivel \texttt{hardware/adc} y \texttt{hardware/dma} del Pico SDK para digitalizar los ejes analógicos del mando (joystick y potenciómetro) con una resolución de 12 bits. Configuré el convertidor analógico-digital en modo \textit{round-robin} sobre cuatro canales a una tasa de $4{,}000$ muestras por segundo por canal ($16{,}000$ conversiones por segundo en total con divisor de reloj de $2999$ a $48\text{ MHz}$), logrando un tiempo de actualización de $0.25\text{ ms}$, cuarenta veces más rápido que el requisito mínimo de $10\text{ ms}$.

Para evitar cualquier sobrecarga sobre el procesador principal, enlacé dos canales DMA en configuración de anillo circular de 8 bytes (\texttt{DMA\_SIZE\_16}), volcando cada muestra directamente hacia el arreglo \texttt{g\_controls.axis} en SRAM. Dado que el DMA en anillo del RP2040 exige alineación en memoria a potencia de dos, aseguré la estructura compartida con \texttt{\_\_attribute\_\_((aligned(8)))} e implementé la función \texttt{read\_adc\_is\_alive()} para supervisar el avance continuo de las conversiones desde el \textit{watchdog}. También colaboré en el armado físico del prototipo y en la documentación de las Clases 8, 9, 10, 15, 18, 20 y 21.

\section*{Aprendizajes Obtenidos}

El diseño de la adquisición analógica por DMA me permitió comprender a fondo las restricciones físicas de los controladores de acceso directo a memoria, tales como la alineación estricta de direcciones y el uso de transferencias atómicas de 16 bits para prevenir lecturas parciales sin necesidad de bloqueos por software.

Asimismo, aplicar una política de asignación de memoria completamente estática sin uso de memoria dinámica (\textit{heap}) demostró cómo se elimina la fragmentación y se garantiza un comportamiento determinista en microcontroladores de recursos limitados.

\vspace{0.3cm}
\noindent\textbf{Repositorio del proyecto:} \\
\url{https://github.com/YizusDP/Practicas_Embebidos/tree/main/ProyectoFinal}

\end{document}